\documentclass[11pt]{article}

% ============================================================
%  COURSEWORK TRANSCRIPT, STYLED AS AN ARXIV PREPRINT
%  Your academic record, submitted to the ledger of history.
% ============================================================

\usepackage[margin=1in]{geometry}
\usepackage[T1]{fontenc}
\usepackage{amsmath, amssymb, amsthm}
\usepackage{graphicx}   % for the rotated arXiv stamp
\usepackage{xcolor}
\usepackage{titlesec}
\usepackage{fancyhdr}
\usepackage{lastpage}
\usepackage[colorlinks=true]{hyperref}

% ---------- The red. THE red. ----------
% #B31B1B is the official arXiv wordmark color (Cornell red).
\definecolor{arxivred}{HTML}{B31B1B}

% ---------- arXiv-flavored link colors ----------
\hypersetup{
  linkcolor = {red!55!black},
  citecolor = {green!45!black},
  urlcolor  = {blue!70!black}
}

% ---------- Section headers in arXiv red ----------
\titleformat{\section}
  {\Large\bfseries\color{arxivred}}{\thesection}{0.8em}{}
\titleformat{\subsection}
  {\large\bfseries\color{arxivred}}{\thesubsection}{0.8em}{}

% ---------- Page style ----------
\pagestyle{fancy}
\fancyhf{}
\renewcommand{\headrulewidth}{0pt}
\fancyfoot[C]{\thepage}

% First page carries the vertical arXiv identifier stamp,
% just like the real thing. math.HO = History and Overview,
% the honest category for a transcript.
\fancypagestyle{firstpage}{%
  \fancyhf{}
  \fancyfoot[C]{\thepage}
  \renewcommand{\headrulewidth}{0pt}
  \fancyhead[L]{\begin{picture}(0,0)
    \put(-52,-480){\rotatebox{90}{\color{arxivred}\small
      \texttt{arXiv:2607.00000v1\ \ [math.HO]\ \ 9 Jul 2026}}}
  \end{picture}}
}

% ---------- Course entries as theorem-style blocks ----------
% Courses are numbered within sections, like results in a paper:
% "Course 2.3" is the third course of the second institution.
\newcounter{course}[section]
\renewcommand{\thecourse}{\thesection.\arabic{course}}

% \course{Grade}{Number}{Title}{Term}{Instructor}
\newcommand{\course}[5]{%
  \refstepcounter{course}%
  \medskip
  \noindent{\color{arxivred}\textbf{Course \thecourse}}\ \textbf{(#2, #4).}\enspace
  \textit{#3.}\enspace Taught by #5.\hfill
  {\color{arxivred}\textbf{[#1]}}\par\nopagebreak
  \smallskip
}
\newcommand{\coursedesc}[1]{\noindent #1\par\nopagebreak}
\newcommand{\coursetext}[1]{\noindent\emph{Textbook:} #1\par}

% Optional: a "remark" for asides (fun facts, war stories)
\newcommand{\remark}[1]{\smallskip\noindent\emph{Remark.}\enspace #1\par}

% ---------- Title block ----------
\title{\bfseries Transcript of Coursework}
\author{Benjamin Wang\\[2pt] 
  \small Department of Mathematics, Dartmouth College\\
  \small Department of Computer Science, Dartmouth College\\
  \small \texttt{benjamin.wang.29@dartmouth.edu}\ \ $\cdot$\ \
  \texttt{\href{https://your-website.example}{https://benjaminw2025.github.io/}}}
\date{July 9, 2026}

% ============================================================
\begin{document}

\maketitle
\thispagestyle{firstpage}

\begin{abstract}
This document presents an overview of my studies in mathematics and computer science coursework conducted at Dartmouth College, with the addition of any self-guided and independent studies at the bottom. Only major related courses will be present (mathematics and computer science). \bf{Grades marked with an * indicate currently enrolled courses}.
\end{abstract}

\noindent\textcolor{arxivred}{\textbf{Keywords:}} probability, compressed sensing, machine learning,
optimization.\\

% ------------------------------------------------------------

\section{Mathematics}

\course{*}{MATH 71}{Algebra}{Fall 2026}{Prof. Ina Petkova}
\coursedesc{Honors track introduction course to abstract algebra. Groups, Group actions, Lagrange's Theorem, Isomorphism Theorems, Sylow Theorems, Rings, Ideals, Quotient rings, Principal Ideal Domains, Unique Factorization Domains.}

\course{A}{MATH 60}{Honors Probability Theory}{Spring 2026}{Instructor Linh Huynh}
\coursedesc{Honors track first year probability theory course. Random variables, order statistics, convolutions, properties of expectation, moment generating functions, limit theorems, concentration inequalities, continuous-time Markov chains, stochastic optimization.}
\coursetext{Ross, \textit{A First Year Course in Probability}; lecture notes.}

\course{A-}{MATH 24}{Linear Algebra}{Winter 2026}{Prof. Dana Williams}
\coursedesc{Dartmouth's equivalent of an abstract linear algebra class; proof variant of MATH 22 Linear Algebra with Applications. Vector spaces, linear operations, matrices, elementary matrix operations, determinants, spectral theory, inner product spaces, SVD.}
\coursetext{Friedberg, \textit{Linear Algebra}; lecture notes.}

\course{A}{MATH 11}{Accelerated Multivariable Calculus}{Fall 2025}{Instructor Gayee Park}
\coursedesc{A fast paced introductory course covering MATH 9 and MATH 13 in one quarter. Vector geometry, space curves, limits and continuity, partial derivatives, tangent planes and differentials, multivariable chain rule, optimization problems, multiple integrals, Green's Theorem, Stoke's Theorem, Divergence Theorem.}
\coursetext{None.}

\section{Computer Science}

\course{*}{COSC 61}{Database Systems}{Fall 2026}{Instructor J. Peter Brady}
\coursedesc{This course introduces the principles behind how modern databases store, organize, query, and reliably manage data, with an emphasis on relational databases and transaction processing. SQL, Relational algebra, Database design, Entity-Relationship (ER) modeling, Normalization, Joins, Transactions, ACID properties, Indexing, Query optimization, Concurrency control, NoSQL/MongoDB.}

\course{A}{COSC 50}{Software Design \& Implementation}{Winter 2026}{Instructor J. Peter Brady}
\coursedesc{An introductory systems programming course injected with lessons on software implementation. Covers Linux, bash commands, memory management, pointers, programming in C. Projects in C include Tiny Search Engine (a miniature web crawler), manually coding fundamental data structures, and an interactive server-client game.}
\coursetext{None}

\course{*}{COSC 39}{Theory of Computation}{Fall 2025}{Prof. Hsieh-Chih Chang}
\coursedesc{This course is an introduction to formal models of languages and computation. Topics covered include finite automata and regular languages, context-free languages, Turing machines and computability, NP-completeness and glimpses of computational complexity theory.}

\course{A}{COSC 30}{Discrete Mathematics for Computer Science}{Winter 2026}{Prof. Hsieh-Chih Chang}
\coursedesc{A discrete math class covering proof techniques (contradiction, induction, pigeonhole), counting, linearity of expectation, stars and bars, strong induction, graph proofs. Very competition math flavored course, with many interesting and fun problems!}
\coursetext{None}

\course{A}{COSC 10}{Problem Solving with Object Oriented Programming}{Fall 2025}{Prof. Alberto Quatrini Li}
\coursedesc{Now renamed to Data Structures and Reasoning. Course covers all major data structures/data types: arrays, linked lists, stacks, queues, binary trees, balanced trees, priority queues, hash tables, graphs; and major graph traversal algorithms: depth-first search, breadth-first search, Dijkstra's algorithm, A* search. In addition, miscellaneous algorithms/topics like Booyer Moore string matching, server client, concurrency and threads, and more.}
\coursetext{Goodrich, \textit{Data Structures and Algorithms}.}

\section{Crosslisted Computer Science}

\course{A}{ENGS 109}{High-Dimensional Sensing and Learning}{Spring 2026}{Prof. Sang Chin}
\coursedesc{A graduate level compressed sensing course. Begins with sparse recovery theory: orthogonal matching pursuit, CoSaMP, subspace pursuit, restricted isometry property, mutual coherence, DCT and random sensing matrices, $\ell_1$ minimization and convex relaxation. Ended with sparse matrix recovery, RPCA, proximal operators, non-negative matrix factorization, dictionary learning algorithms, IHT, LISTA. Research styled problem sets, and a poster presentation in addition to a written final.}
\coursetext{Foucart and Rauhut, \textit{A Mathematical Introduction to
Compressive Sensing}.}
\remark{Implemented every matching pursuit algorithm from scratch before reading the reference
implementation. Recommended.}

\course{*}{ENGS 101}{Principles of Reinforcement Learning}{Fall 2026}{Prof. Sang Chin}
\coursedesc{A graduate level course on reinforcement learning. Covers all major mathematical foundations, algorithms, and applications for reinforcement learning.}
\end{document}